\section{Fixed shifts and wide major arcs}\label{sec:fixed-shift-comparison}

For fixed \(k=h\ne0\), the shared-label comparison
\eqref{eq:gr-comparison}--\eqref{eq:gr-comparison-term} remains valid.
The frequency estimate has a different scale. After localizing
\(D_i\asymp H_i\), put
\[
 H'=H_1H_2,\qquad Y=xH',\qquad H=H'L^{-B_0}.
\]
The multiplier is now \(m(h\theta)\). The integer bilinear bound
\eqref{eq:gr-multiplier} confines its significant part to
\(O_h(L^{2B_0})\) arcs about rationals \(a/q\), with
\[
 q\le |h|L^{B_0},\qquad
       |\theta-a/q|\le C_h/H.
\]
In particular, their physical width is \(L^{B_0}/H'\), rather than
\(L^{B_0}/Y\). The target-grid proposition does not estimate these
arcs: on them \(|\beta|Y\) can be of order \(xL^{B_0}\).
We extract the already established Dirichlet-frequency argument of
\cite[Sections 6.2--6.4]{walsh}, with the arithmetic hypotheses
\textup{(C1)--(C3)} in place of its particular divisor model.

\begin{proposition}[Fixed-shift comparison energy]
\label{prop:fixed-comparison-energy}
Let \(F\) have the form \eqref{eq:gr-energy-F}, with a selected
small bank whose label does not enter \(b\), and with its exact
prime count equal to one. Suppose
\[
 H_i\ge\exp(L^{1/3}),\qquad
 H_i\le\exp(O((\log L)L^{2/5})),\qquad Y=xH_1H_2.
\]
For every fixed \(B_0,D>0\), conditions \textup{(C1)--(C3)} and
the zero-free input imply, at sufficiently deep fixed model depth,
\begin{equation}\label{eq:fixed-wide-energy}
 \int_{|\beta|\le C_h/H}
      |\widehat F(a/q+\beta)|^2\,d\beta
                            \ll_{D,h}YL^{-D}
\end{equation}
uniformly for \(q\le |h|L^{B_0}\), where
\(\widehat F(\theta)=\sum_nF(n)\ee{\theta n}\).
Consequently every fixed-shift comparison in
\eqref{eq:gr-comparison-term} has an arbitrary prescribed
logarithmic saving. No binary cancellation estimate is assumed.
\end{proposition}

\subsection{Conversion to Dirichlet frequencies}

We use the following existing two mean-square inequalities, in their
normalizations from \cite[Lemma 6.3 and (6.14)]{walsh}. If \(a_n\)
is supported on \(n\asymp Y\), then
\begin{equation}\label{eq:fixed-additive-mellin}
 \int_{|\beta|\le C/H}\left|\sum_na_n\ee{\beta n}\right|^2d\beta
 \ll_C\frac1Y\int_{\mathbb R}(1+H|t|/Y)^{-2}
                         \left|\sum_na_nn^{it}\right|^2dt.
\end{equation}
Here is the averaging step in that inequality, following
\cite[Lemmas 2.1--2.2]{soundararajan}. Put
\(\eta=H/(4C)\), \(A(v)=\sum_{n\le v}a_n\),
\(P(t)=\sum_na_nn^{it}\), and
\(J_\eta=\int_0^\infty|A(v+\eta)-A(v)|^2dv\).
The interval kernel has Fourier transform of modulus at least
\(c\eta\) on \(|\beta|\le C/H\), so Plancherel bounds the
left side by \(O(\eta^{-2}J_\eta)\).
For \(2\eta\le\nu\le3\eta\), split the increment at \(v+\nu\),
square, and average in \(\nu\). Changing \(v\) to \(v+\eta\)
in the second increment gives, if the coefficients are supported
on \([aY,bY]\),
\[
 J_\eta\ll\eta^{-1}\int_\eta^{3\eta}
       \int_{aY/2}^{2bY}|A(v+\nu)-A(v)|^2dv\,d\nu.
\]
Substitute \(\nu=\Delta v\). Since \(v\asymp Y\), this is at most
\[
 \frac{CY^2}{\eta}
 \int_{c_1\eta/Y}^{c_2\eta/Y}\int_0^\infty
       |A(v(1+\Delta))-A(v)|^2\frac{dv}{v}\,d\Delta.
\]
Logarithmic Plancherel evaluates the inner integral exactly as
\[
 \frac1{2\pi}\int_{\mathbb R}|P(t)|^2
          \left|\frac{(1+\Delta)^{it}-1}{t}\right|^2dt.
\]
The kernel is \(O(\min\{\eta^2/Y^2,|t|^{-2}\})\), and
the \(\Delta\)-interval has length \(O(\eta/Y)\). Thus
\(J_\eta\ll Y\int\min\{\eta^2/Y^2,|t|^{-2}\}|P(t)|^2dt\).
Multiplying by \(\eta^{-2}\) proves
\eqref{eq:fixed-additive-mellin}; no identification of fixed additive
windows with fixed multiplicative windows is used. The other inequality,
the ordinary Dirichlet-polynomial mean-value theorem, is
\begin{equation}\label{eq:fixed-dirichlet-mean}
 \int_I\left|\sum_{n\le CR}a_nn^{it}\right|^2dt
              \ll_C (|I|+R)\sum_n|a_n|^2,
\end{equation}
uniformly in the location of the interval \(I\). Neither inequality
is a new distributional input.

Apply \eqref{eq:fixed-additive-mellin} to
\(a_n=F(n)\ee{an/q}\) and set
\[
 \mathcal T(t)=Y^{-1}\sum_nF(n)\ee{an/q}n^{it},\qquad
 T_0=(Y/H)L^{B_1}=xL^{B_0+B_1}.
\]
The pointwise list bound and \textup{(C1)} give
\(\sum_n|F(n)|^2\ll YL^C\). Thus on a dyadic interval
\(T\le |t|\le2T\), \eqref{eq:fixed-dirichlet-mean} gives
\[
 \int_{T\le|t|\le2T}|\mathcal T(t)|^2dt
                                  \ll L^C(1+T/Y).
\]
Summing this with the weight in
\eqref{eq:fixed-additive-mellin}, for \(|t|>T_0\), gives
\begin{equation}\label{eq:fixed-frequency-tail}
 \int_{|t|>T_0}(1+H|t|/Y)^{-2}|\mathcal T(t)|^2dt
       \ll L^C\bigl(L^{-2B_1}+H^{-1}L^{-B_1}\bigr).
\end{equation}
The exponent \(B_1\) is chosen after the required saving, so this
tail is arbitrarily small. It remains to prove
\begin{equation}\label{eq:fixed-dirichlet-energy}
                \int_{|t|\le T_0}|\mathcal T(t)|^2dt\ll_D L^{-D}.
\end{equation}

\subsection{Low frequencies and recordwise centering}

Choose a fixed exponent \(B_2\), to be specified after the smooth
prime estimates, and put \(T_c=L^{B_2}\). In the low range write
\(n=rv\), where \(r=n_S\) is the complete smooth record and
\(v\) is rough. Every list weight is a function of \(r\) alone
and is bounded by \(L^C\). The terms with \(r>E\) have population
\(O(YL^C\exp(-cL^{1/10}))\), by the Rankin estimate in
\eqref{eq:gr-tailsets}; their contribution to \(\mathcal T(t)\)
is negligible uniformly in \(t\). For each remaining record,
\(\ee{arv/q}\) is a bounded periodic test of period at most \(q\).
The phase \((rv)^{i(t+\sigma)}\psi(rv/Y)\) has fixed logarithmic
derivatives bounded by fixed powers of \(L\), throughout
\(|t|\le T_c\). Condition \textup{(C3)} therefore gives any
prescribed saving for that record. Summing uses only
\[
          \sum_{r\ S\text{-smooth}}\frac1r
             =\prod_{p\le y}(1-1/p)^{-1}\ll L^C.
\]
Choose the precision in \textup{(C3)} after \(B_2\) and all list
costs. This gives
\begin{equation}\label{eq:fixed-low}
             \int_{|t|\le T_c}|\mathcal T(t)|^2dt\ll_D L^{-D}.
\end{equation}
The arithmetic progression test is on \(v\), so this argument
does not require \(r\) or the endpoint to be a unit at \(L_0\).

\subsection{The smooth prime estimate at polynomial height}

The high-frequency argument uses a smooth version of
Lemma~\ref{lem:gr-primes}. Its height dependence is essential:
the linear \(1+|t|\) loss in the sharp-interval estimate is not
adequate when \(|t|\) can be of order \(x\).
Fix \(\sigma_1\in(11/12,1)\). For \(w\) smooth and supported on
\([1/2,2]\), with each required derivative bounded by a fixed
power of \(L\), the same zero-free input gives
\begin{equation}\label{eq:fixed-smooth-primes}
 \sum_p\chi(p)w(p/P)p^{it}
  =\delta_\chi\int_0^\infty\frac{w(u/P)u^{it}}{\log u}\,du
                       +O(P^{\sigma_1}L^{C_w})
\end{equation}
uniformly for \(q\le L^B\), characters \(\chi\bmod q\),
\(|t|\le x^C\), and \(P\ge\exp(cL^{1/8})\).
Here \(C_w\) depends on the fixed derivative bounds and on \(B,C\).
This is the smooth prime estimate used in
\cite[Lemma 5.1 and Section 6.4]{walsh}.

To obtain it directly from the present input, insert the von
Mangoldt weight and the smooth factor \(w(u/P)/\log u\), apply
Mellin inversion, and move the line from \(\Re s>1\) to
\(\Re s=\sigma_1\). Formula \eqref{eq:gr-logderivative} bounds
the regularized logarithmic derivative on this line by
\(O(\log(q(|\Im s|+3)))\). The Mellin transform is translated
in height by \(t\); its decay depends on the derivatives of
\(w\), rather than on \(|t|\). Its weighted integral is consequently
\(O(L^{C_w})\) even for \(|t|\le x^C\). The principal pole
gives the integral in \eqref{eq:fixed-smooth-primes}; prime powers
cost \(O(P^{1/2}L^{C_w})\), absorbed by the displayed error.
Truncation followed by dominated convergence justifies the contour
move. This uses the same fixed zero-free half-plane, without a
second prime-distribution theorem.

Impose simultaneously the bank class \(c\pmod{L_0}\) and a
prime class modulo \(q\). Incompatible classes contribute zero;
otherwise character orthogonality at \(\lcm(q,L_0)\) applies.
For large \(x\), every prime in the small bank exceeds
\(qL_0|h|\), so nonreduced prime classes and primes dividing
\(h\) do not occur. Integration by parts in the principal
integral, for every fixed \(j\), shows that the normalized
prime sum \(P^{-1}\sum_p\cdots\) is bounded by
\begin{equation}\label{eq:fixed-high-prime}
        L^{C_{w,j}}(1+|t|)^{-j}
                       +L^{C_w}\exp(-cL^{1/8}).
\end{equation}
All character and congruence costs here are polynomial logarithmic.

\subsection{High frequencies and a long cofactor}

For this part remove selected-bank prime squares, whose population
is \(O(Y/P_{\min})\). Applied to their bounded coefficients,
\eqref{eq:fixed-dirichlet-mean} bounds their normalized integrated
energy on \([-T_0,T_0]\) by
\(L^C(1+T_0/Y)/P_{\min}\), hence by less than every logarithmic
power. Do not impose the cutoff \(n_S\le E\) here: it would
couple the selected prime to its cofactor.

The remaining terms have a unique selected-bank prime and can be
written \(n=pm\). Exact invariance \textup{(C2)} gives
\(\xi(pm)=\xi(m)\). The distinguished list coefficient ignores
this prime, so after removing its factor \(V_g^{-1}\) the
remaining coefficient \(\gamma(m)\) is independent of \(p\) and
bounded by \(L^C\). It contains the condition that \(m\) has no
selected-bank prime, together with the other list weights.

Partition the selected bank into smooth dyadic pieces
\(p\asymp P\), including its two sharp outer endpoints. Replace
each sharp endpoint by a smooth cutoff across a strip of length
\(\varepsilon P\), where \(\varepsilon=L^{-A_{\rm sm}}\).
The normalized sum of prime labels in such a strip is at most
\(O(\varepsilon+P^{-1})\), simply by counting integers. The
Dirichlet mean-value bound for the cofactor below and Minkowski
show that this replacement changes the normalized integrated
energy by at most
\[
                  L^C(\varepsilon+P_{\min}^{-1})^2.
\]
There are only polynomially many dyadic pieces in \(L\).
Choose \(A_{\rm sm}\) large after the desired saving and these
costs. The resulting smooth prime cutoffs have derivative bounds
which are fixed powers of \(L\).

Put \(R=Y/P\), and insert a smooth cofactor cutoff
\(g(m/R)\) equal to one wherever the product cutoff can contribute.
Separate the additive phase by the residue classes of \(p,m\)
modulo \(q\), and Mellin-invert the product cutoff. Apart from
the errors just bounded, \(\mathcal T(t)\) is a sum of at most
polynomially many terms of the form
\begin{equation}\label{eq:fixed-prime-cofactor}
 \frac{1}{2\pi V_g}\sum_{r,r'\bmod q}\ee{arr'/q}
       \int_{\mathbb R}\widetilde\psi(\nu)Y^{-i\nu}
          P_r(t+\sigma+\nu)R_{r'}(t+\sigma+\nu)\,d\nu,
\end{equation}
where harmless Mellin sign choices are fixed by this formula and
\[
 \begin{split}
 P_r(s)&=P^{-1}\sum_{\substack{p\equiv r\ (q)\\p\equiv c\ (L_0)}}
                    w(p/P)p^{is},\\
 R_{r'}(s)&=R^{-1}\sum_{m\equiv r'\ (q)}
                      g(m/R)\gamma(m)m^{is}.
 \end{split}
\]
There is no unit restriction on \(r'\). For any interval of
length at most \(4T_0\), \eqref{eq:fixed-dirichlet-mean} gives,
uniformly in its location,
\begin{equation}\label{eq:fixed-long-cofactor}
       \int_I|R_{r'}(t)|^2dt\ll L^C(1+T_0/R)\ll L^C.
\end{equation}
The final inequality uses the decisive scale separation
\[
 \frac{T_0}{R}\ll\frac{P L^{B_0+B_1}}{H_1H_2}
   \le\exp\bigl(CL^{1/4}-2L^{1/3}\bigr)L^{B_0+B_1}=o(1).
\]
It comes from the two independent big-prime products already
present in the graph comparison, rather than a new energy bound.

The Mellin transform in \eqref{eq:fixed-prime-cofactor} has
Schwartz decay with polynomial-logarithmic seminorms. Choose a
fixed exponent \(C_\nu\) large enough that its tail
\(|\nu|>L^{C_\nu}\), in weighted \(L^1\), is smaller than the
required logarithmic power. Uniformity in the location of \(I\)
in \eqref{eq:fixed-long-cofactor}, and the trivial bound
\(|P_r|\ll1\), justify removing this tail. Now choose \(B_2\)
larger than \(C_\nu\), the bound for \(|\sigma|\), and the
smooth-cutoff derivative costs. Throughout
\(T_c<|t|\le T_0\), the argument
\(s=t+\sigma+\nu\) then has \(|s|\gg L^{B_2}\).
In \eqref{eq:fixed-high-prime} fix, for example, the
integration-by-parts order \(j=2\) first. After choosing the
requested exponent \(A\), enlarge \(B_2\) beyond the resulting
fixed derivative costs. This gives
\[
           \sup_{T_c<|t|\le T_0,\ |\nu|\le L^{C_\nu}}
                    |P_r(t+\sigma+\nu)|\ll_A L^{-A}
\]
for any requested \(A\). The smooth prime estimate applies
because \(T_0=xL^{B_0+B_1}\le x^2\) for large \(x\).
Choose this \(A\) after the fixed partition and Mellin costs.
Combine the estimate with \eqref{eq:fixed-long-cofactor}, then apply
Minkowski to the Mellin integral and sum the dyadic and residue
pieces. This proves
\[
              \int_{T_c<|t|\le T_0}|\mathcal T(t)|^2dt\ll_D L^{-D}.
\]
Together with \eqref{eq:fixed-low} this is
\eqref{eq:fixed-dirichlet-energy}; together with
\eqref{eq:fixed-frequency-tail} and
\eqref{eq:fixed-additive-mellin} it proves
\eqref{eq:fixed-wide-energy}.

The parameter order is explicit. First fix the graph comparison,
its list bounds, and its external cutoff costs. Then fix the major
arcs and the required energy saving, the tail cutoff \(B_1\),
the bank smoothing precision, and the Mellin truncation. Choose
the high-frequency boundary \(B_2\) after those choices. Finally
choose the model depth through \textup{(C3)} at the required
low-frequency precision. No exponent in the graph or the endpoint
moment bound depends on that final depth.

On the arcs, Cauchy--Schwarz, \eqref{eq:fixed-wide-energy}, and
Parseval for the other endpoint bound the normalized Fourier
comparison by every requested inverse power of \(L\), after
paying the polynomially many arcs. Off the arcs,
\eqref{eq:gr-multiplier} and the two endpoint norms do so.
The normalization is \(1/Y\), exactly as in
\eqref{eq:gr-fouriercomparison}, with \(T\) replaced by \(h\).
The harmonic conversion \(\Psi(t)=t\Phi(t)\) therefore gives
the same centered banked replacement as for reflected sums.
This completes the proof of
Proposition~\ref{prop:fixed-comparison-energy} and supplies the
fixed-shift branch used in Section~\ref{sec:assembly}.
